\subsection{并集与交集的像}

命题3。设 \((\mathbf{X}_t)_{t\in I}\) 是集合A的子集族，并设 \(\Gamma\) 是A与B之间的对应关系。则

\[
\Gamma \left\langle \bigcup_ {i \in I} X _ {i} \right\rangle = \bigcup_ {i \in I} \Gamma \left\langle X _ {i} \right\rangle , \quad \Gamma \left\langle \bigcap_ {i \in I} X _ {i} \right\rangle \subset \bigcap_ {i \in I} \Gamma \left\langle X _ {i} \right\rangle .
\]

关系 \((\exists x)\left(x\in \bigcup_{i\in I}X_i \text{ 且 } y\in \Gamma (x)\right)\) 等价于

\[
(\exists x) (\exists t) (t \in I \text {   且   } x \in X _ {t} \text {   且   } y \in \Gamma (x)),
\]

因此到 \((\exists \iota)(\iota \in I \text{ 且 } y \in \Gamma\langle X_{\iota}\rangle)\) ，因此等价于 \(y \in \bigcup_{\iota \in I} \Gamma\langle X_{\iota}\rangle\) ；这证明了第一个公式。至于第二个公式，我们有 \(\bigcap_{\iota \in I} X_{\iota} \subset X_{\iota}\) 对于所有 \(\iota \in I\) ，由此（§3，命题2）

\[
\Gamma \left\langle \bigcap_ {i \in I} X _ {i} \right\rangle \subset \Gamma \left\langle X _ {i} \right\rangle ,
\]

，从而

\[
\Gamma \left\langle \bigcap_ {i \in I} X _ {i} \right\rangle \subset \bigcap_ {i \in I} \Gamma \left\langle X _ {i} \right\rangle .
\]

如果 \(\Gamma\) 是任意对应关系（特别是任意函数），公式

\[
\Gamma \left\langle \bigcap_ {i \in I} X _ {i} \right\rangle = \bigcap_ {i \in I} \Gamma \left\langle X _ {i} \right\rangle
\]

通常不成立。

* 例如，在平面中 \(\mathbf{R}^2\) 这些直线的第一投影 \(y = x\) 并且 \(y = x + 1\) 等于 \(\mathbf{R}\) , 但这些直线的交集为空，因此该交集的第一投影也为空 (*)。*

然而，我们有以下重要结果：

命题4。设 \(f\) 设 为从 A 到 B 的映射，并设 \((\mathbf{Y}_t)_{t \in I}\) 为 B 的子集族。则 \(\overline{f}^1 \left\langle \bigcap_{i \in I} \mathbf{Y}_t \right\rangle = \bigcap_{i \in I} \overline{f}^1 \langle \mathbf{Y}_t \rangle\) .

为证明这一点，设 \(x\) 为 \(\bigcap_{i \in I} f^{-1} \langle Y_i \rangle\) 我们有 \(f(x) \in Y_i\) 对于所有 \(i \in I\) ，由此 \(f(x) \in \bigcap_{i \in I} Y_i\) ，并且因此

\[
x \in \overline {{f}} ^ {1} \left\langle \bigcap_ {\iota \in I} Y _ {\iota} \right\rangle .
\]

因此 \(\bigcap_{i\in I}\overline{f}^{1}(Y_{i})\subset\overline{f}^{1}\Big\langle\bigcap_{i\in I}Y_{i}\Big\rangle\) ；这一关系，连同命题3，给出了结果。

推论. 如果 \(f\) 是从 \(A\) 到 \(B\) 的单射，且 \((X_{t})_{t \in I}\) 是 \(A\) 的子集族，其指标集 \(I\) 非空，则 \(f\left\langle \bigcap_{i \in I} X_{t} \right\rangle = \bigcap_{i \in I} f\langle X_{t} \rangle\) .

因为可写成 \(f = i \circ g\) ，其中 \(i\) 是从 \(f\langle \mathbf{A} \rangle\) 到 \(\mathbf{B}\) 的典范单射，而 \(g\) 是从 \(\mathbf{A}\) 到 \(f\langle \mathbf{A} \rangle\) 的双射。如果 \(h\) 表示 \(g\) 的逆映射，则我们有 \(f\langle \mathbf{X} \rangle = \overline{h}^1\langle \mathbf{X} \rangle\) ，其中 \(\mathbf{X}\) 是 \(\mathbf{A}\) 的任意子集，因而归结为命题4。

